\documentclass[a4paper, serif, 10pt]{moderncv}

%% moderncv
% color and style
\moderncvcolor{black}
\moderncvstyle{banking}

% renew moderncv command to adapt for CJK colon
\renewcommand*{\cvitem}[3][.25em]{%
  \ifstrempty{#2}{}{\hintstyle{#2}：}{#3}%
  \par\addvspace{#1}}

\renewcommand*{\cvitemwithcomment}[4][.25em]{%
  \savebox{\cvitemwithcommentmainbox}{\ifstrempty{#2}{}{\hintstyle{#2}：}#3}%
  \setlength{\cvitemwithcommentmainlength}{\widthof{\usebox{\cvitemwithcommentmainbox}}}%
  \setlength{\cvitemwithcommentcommentlength}{\maincolumnwidth-\separatorcolumnwidth-\cvitemwithcommentmainlength}%
  \begin{minipage}[t]{\cvitemwithcommentmainlength}\usebox{\cvitemwithcommentmainbox}\end{minipage}%
  \hfill% fill of \separatorcolumnwidth
  \begin{minipage}[t]{\cvitemwithcommentcommentlength}\raggedleft\small\itshape#4\end{minipage}%
  \par\addvspace{#1}}

%% page layout/margins
\usepackage[top=2.5cm, bottom=2.5cm, left=1.5cm, right=1.5cm]{geometry}

\nopagenumbers{}

%% Babel config for English language
\usepackage[english]{babel}

%% fontspec
\usepackage{fontspec}

\IfFontExistsTF{Linux Libertine}{
  \setmainfont[Ligatures={TeX, Common}, Numbers=Lining, ItalicFont=Linux Libertine]{Linux Libertine}
}{}
\IfFontExistsTF{Linux Libertine O}{
  \setmainfont[Ligatures={TeX, Common}, Numbers=Lining, ItalicFont=Linux Libertine O]{Linux Libertine O}
}{}

%% CTeX
% CJK support, used to show CJK characters in the resume
%
% - heading=false: disable ctex heading
% - punct=kaiming: use kaiming punctuations styles for CJK
% - scheme=plain: use plain scheme, do not override \normalsize font size
% - space=auto: space settings for CJK characters
%
% CJK fonts are resolved per locale language, e.g. Noto Serif/Sans CJK JP
% for Japanese resumes, with fallback to the next region if not installed
%
% ref:
% - http://ctan.mirrorcatalogs.com/language/chinese/ctex/ctex.pdf
\usepackage[UTF8, heading=false, punct=kaiming, scheme=plain, space=auto]{ctex}

\IfFontExistsTF{Noto Serif CJK JP}{
  \setCJKmainfont{Noto Serif CJK JP}
}{}
\IfFontExistsTF{Noto Sans CJK JP}{
  \setCJKsansfont{Noto Sans CJK JP}
}{}

%% line spacing
\usepackage{setspace}
\setstretch{1.125}

%% URL styling - use normal text instead of monospace
\urlstyle{same}

%% Patch \httplink and \httpslink to handle full URLs
% The original moderncv macros always prepend http:// or https:// to URLs.
% This causes issues when the URL already has a protocol (e.g., https://example.com)
% resulting in malformed URLs like https://https://example.com.
% This patch checks if the URL already contains "://" and uses it directly if so.
% Note: We use \str_set:Nx (x-expansion) to expand macros like \@homepage first.
\ExplSyntaxOn
\str_new:N \g_yamlresume_url_str

\RenewDocumentCommand{\httpslink}{O{}m}{
  \str_set:Nx \g_yamlresume_url_str {#2}
  \str_if_in:NnTF \g_yamlresume_url_str {://}
    { \url{#2} }
    { \url{https://#2} }
}

\RenewDocumentCommand{\httplink}{O{}m}{
  \str_set:Nx \g_yamlresume_url_str {#2}
  \str_if_in:NnTF \g_yamlresume_url_str {://}
    { \url{#2} }
    { \url{http://#2} }
}
\ExplSyntaxOff

\name{佐藤 健一}{}
\title{サイバーセキュリティエンジニア}
\phone[mobile]{+81 90-1234-5678}
\email{kenichi.sato@example.com}
\homepage{https://kenichi-sato.dev}

\address{日本、東京都、千代田区、東京都千代田区丸の内1-2-3、100-0001}{}{}

\extrainfo{{\small \faGithub}\ \href{https://github.com/kenichi-sato}{@kenichi-sato} {} {} {} • {} {} {} 
{\small \faLinkedin}\ \href{https://www.linkedin.com/in/kenichi-sato}{@kenichi-sato} {} {} {} • {} {} {} 
{\small \faTwitter}\ \href{https://twitter.com/kenichi\_sato}{@kenichi\_sato}}

\begin{document}

\maketitle

\section{基本情報}

\cvline{}{セキュリティエンジニアとして5年以上の経験を持ち、ネットワークセキュリティ、脆弱性診断、インシデント対応に精通しています。
特にペネトレーションテストとセキュリティアーキテクチャの設計に強みを持ち、金融機関やEコマース企業のセキュリティ強化に貢献してきました。
チームをリードし、セキュリティ意識向上のためのトレーニングを実施した経験もあります。
常に最新の脅威情報を追い、プロアクティブな防御策を提案することを心がけています。}
\section{学歴}

\cventry{2013年4月～2015年3月}
        {修士、情報理工学系研究科 コンピュータ科学専攻、成績：3.8/4.0}
        {東京大学}
        {\url{https://www.u-tokyo.ac.jp}}
        {}
        {情報セキュリティを専攻し、暗号プロトコルとネットワーク防御に関する研究に従事。
修士論文では「機械学習を用いた異常検知による侵入検知システムの改善」をテーマに研究。
\textbf{コース}：情報セキュリティ、暗号理論、ネットワーク工学、オペレーティングシステム、データベースシステム}
\section{職歴}

\cventry{2019年4月～現在}
        {シニアセキュリティエンジニア}
        {株式会社サイバーディフェンス}
        {\url{https://www.cyberdefense.co.jp}}
        {}
        {\begin{itemize}
\item 金融機関向けのセキュリティアーキテクチャ設計と実装をリード
\item 脆弱性診断およびペネトレーションテストを実施し、重大な脆弱性を複数発見・修正
\item インシデント対応チームのリーダーとして、セキュリティインシデントの調査と復旧を指揮
\item セキュリティポリシーの策定と従業員向けセキュリティトレーニングを実施
\item 最新の脅威情報を収集し、プロアクティブな防御策を提案
\end{itemize}
\textbf{キーワード}：ネットワークセキュリティ、ペネトレーションテスト、インシデント対応、セキュリティアーキテクチャ}

\cventry{2015年4月～2019年3月}
        {セキュリティエンジニア}
        {株式会社セキュアテック}
        {\url{https://www.securetech.co.jp}}
        {}
        {\begin{itemize}
\item Webアプリケーションの脆弱性診断を担当し、SQLインジェクションやXSSなどの脆弱性を特定
\item セキュリティ監視システムの運用と改善を実施
\item セキュリティインシデントの一次対応と報告書作成
\item 社内セキュリティガイドラインの整備に貢献
\end{itemize}
\textbf{キーワード}：脆弱性診断、Webセキュリティ、セキュリティ監視}
\section{言語}

\cvline{日本語}{ネイティブ・母国語レベル \hfill \textbf{キーワード}：日本語ネイティブ}
\cvline{英語}{完全な実務レベル \hfill \textbf{キーワード}：TOEIC 900点、技術文書読解}
\section{スキル}

\cvline{ネットワークセキュリティ}{エキスパート \hfill \textbf{キーワード}：ファイアウォール、IDS/IPS、VPN、パケット解析、ネットワーク設計}
\cvline{ペネトレーションテスト}{上級 \hfill \textbf{キーワード}：Metasploit、Burp Suite、Nmap、Wireshark、Kali Linux}
\cvline{クラウドセキュリティ}{中級 \hfill \textbf{キーワード}：AWS、Azure、GCP、IAM、セキュリティグループ}
\cvline{暗号技術}{上級 \hfill \textbf{キーワード}：PKI、TLS/SSL、暗号化アルゴリズム、鍵管理}
\section{受賞・表彰}

\cventry{2022年12月}
        {株式会社サイバーディフェンス}
        {年間最優秀セキュリティエンジニア賞}
        {}
        {}
        {金融機関向けのセキュリティプロジェクトにおける卓越したリーダーシップと技術的貢献に対して授与。}

\cventry{2015年3月}
        {東京大学}
        {セキュリティ研究優秀賞}
        {}
        {}
        {修士論文「機械学習を用いた異常検知による侵入検知システムの改善」が高く評価され授与。}
\section{資格・免状}

\cventry{2020年6月}
        {(ISC)²}
        {CISSP（情報システムセキュリティ専門家）}
        {\url{https://www.isc2.org/Certifications/CISSP}}
        {}
        {}

\cventry{2018年11月}
        {Offensive Security}
        {OSCP（Offensive Security Certified Professional）}
        {\url{https://www.offensive-security.com/pwk-oscp/}}
        {}
        {}

\cventry{2017年4月}
        {情報処理推進機構}
        {情報処理安全確保支援士}
        {\url{https://www.ipa.go.jp/shiken/kubun/sc.html}}
        {}
        {}
\section{出版・著作}

\cventry{2015年3月}
        {機械学習を用いた異常検知による侵入検知システムの改善}
        {情報処理学会 論文誌}
        {\url{https://www.ipsj.or.jp}}
        {}
        {侵入検知システムの精度向上を目的として、機械学習アルゴリズムを適用した異常検知手法を提案。
従来手法と比較して検知率が15\%向上し、誤検知率が20\%低減することを示した。}
\section{プロジェクト}

\cventry{2021年1月～2021年6月}
        {複数のセキュリティログを統合し、リアルタイムで可視化するダッシュボード}
        {セキュリティ監視ダッシュボード}
        {\url{https://github.com/kenichi-sato/security-dashboard}}
        {}
        {\begin{itemize}
\item Elasticsearch、Logstash、Kibana (ELKスタック) を使用してセキュリティログを収集・分析
\item 異常検知アラートをSlackに通知する機能を実装
\item セキュリティ運用チームの効率を30\%向上
\end{itemize}
\textbf{キーワード}：セキュリティ監視、ELK、ログ分析、可視化}
\section{興味・関心}

\cvline{CTF (Capture The Flag)}{Web、暗号、フォレンジック}
\cvline{読書}{セキュリティ関連書籍、技術書}
\section{ボランティア活動}

\cventry{2018年4月～2023年3月}
        {セキュリティトレーナー}
        {日本サイバーセキュリティコミュニティ}
        {\url{https://www.jcsc.org}}
        {}
        {\begin{itemize}
\item 初心者向けのセキュリティワークショップを定期的に開催
\item CTF大会の運営スタッフとして参加
\item セキュリティ意識向上のための講演を実施
\end{itemize}}
    
\end{document}